\chapter{线性代数与多重线性代数。}

线性代数既是数学中最古老的分支之一，又是最新近的分支之一。一方面，在数学的起源处就可找到一些通过单次乘法或除法即可解决的问题，也就是说，通过计算函数 \(f(x) = ax\) 的一个值，或通过解方程 \(ax = b\) 来解决：这些正是线性代数的典型问题，而不“线性地”思考，就不可能处理它们，甚至不可能正确地陈述它们。

另一方面，不仅这些问题，而且几乎一切涉及一次方程的内容，长期以来都被贬入初等教学，直到域、环、拓扑向量空间等概念的现代发展，才解放了线性代数的那些本质概念（例如对偶），并赋予它们重要性；正是在那时，人们认识到了几乎整个现代代数在本质上是线性的——这种“线性化”甚至是其显著特征之一——线性代数也重新获得了它应有的地位。因此，从我们所采取的观点出发来叙述它的历史，将是一项既重要又困难的任务；我们在这里只能满足于相当简略的提示。

由前文所述可知，线性代数无疑是为了满足从事实际计算者的需要而诞生的；因此我们看到，三率法（{[}311{]}，第 I 卷，第~150-155~页）与试位法以或多或少明确的形式，在一切实用算术教科书中扮演着重要角色，从埃及人的莱因德纸草书，经阿耶波多、阿拉伯人、比萨的莱昂纳多以及中世纪与文艺复兴时期不可胜数的“算书”，直到今日我们小学里仍在使用的教科书；但它们或许从来不过是供从业者之用、从更高级的科学理论中摘出的一部分。

至于真正的数学家，他们关于线性代数的研究的性质，乃是其科学的一般结构的函数。古希腊数学，如欧几里得《原本》中所阐述的那样，发展了两个具有线性特征的抽象理论：一方面是量的理论（{[}107{]}，第 V 卷；参见第~152~页），另一方面是整数的理论（{[}107{]}，第 VII 卷）。在巴比伦人那里，我们发现了一些与我们初等代数接近得多的方法；他们知道如何求解一次方程组，而且方式十分优雅（{[}232{]}，第~181-183~页）。然而，在很长一段时期内，线性代数的进步主要是与代数运算的进步联系在一起的，而正是从这种为本注记所不涉及的角度去考察它才是合理的；事实上，要把一个线性方程组化归为 \(ax = b\) 型的方程，当只涉及一个未知数时，只需知道那些（丢番图已经实质上陈述过的）把各项从一个等号边移到另一边并合并同类项的规则就够了；而当涉及多个未知数时，只需再知道如何逐个消去它们，直到只剩下一个即可。因此，直到 18 世纪的那些代数论著都认为，一旦列出了这些规则，在一次方程方面就已经做完了一切；至于未知数个数与方程个数相同的方程组（他们不考虑别的情形），当其各方程的左端不是线性无关的型时，他们总是满足于顺带指出这意味着问题提得不恰当。在 19 世纪的论著中，甚至在某些更晚近的著作中，这一观点也只是靠记号的进步——它使写出 \(n\) 个未知数的 \(n\) 个方程的方程组成为可能——以及行列式的引入——它使在“一般情形”下陈述显式求解的公式成为可能——才有所改变；这一进步，倘若莱布尼茨（{[}198 a{]}，第 I 卷，第~239~页）曾发表过他在这方面的想法，本应归功于他，而它实际上主要归功于 18 世纪及 19 世纪初的数学家们。

但是，我们必须首先考察几股思想潮流，它们对线性代数——按我们所理解的意义——的发展所作的贡献，远比对线性方程组的研究为大。费马受阿波罗尼奥斯著作的启发（{[}109{]}，第 I 卷，第~91-110~页；法文译文，第 III 卷，第~85-101~页），甚至在笛卡儿（{[}85 a{]}，第 VI 卷）之前就构思了解析几何的原理，想到了借助次数对平面曲线进行分类（这一分类到 17 世纪末已为所有数学家越来越熟悉，可以认为已确定地被掌握），并陈述了基本原理：在平面上，一次方程表示一条直线，二次方程表示一条圆锥曲线：他由此立即推出了一些关于几何轨迹的“非常漂亮的”结论。与此同时，他预告了（{[}109{]}，第 I 卷，第~184-188~页；法文译文，第 III 卷，第~159-163~页）对问题的这样一种分类：已解决的问题；可化归为含两个未知数的方程的问题；可化归为含三个未知数的方程的问题；等等；他还补充道：第一类在于确定一个点，第二类在于确定一条直线或一个平面轨迹，其后各类在于确定一个曲面；等等。（“\ldots 这样一个问题所寻求的不仅是一个点或一条线，而是一个与问题相称的整个曲面；正是由此产生了那些面轨迹，其后各类依此类推”，出处同上，第~186~页；在那里已经可以找到 n 维几何的胚胎。）这段陈述了代数中和几何代数中的维数原理的文字，显示出了代数与几何的一种融合，它与现代观念完全合拍，然而如前所见，它花了两个多世纪才深入人心。

至少这些思想不久便导致了解析几何的兴盛，它在 18 世纪随着克莱罗、欧拉、克拉默、拉格朗日以及许多其他人而获得其全部规模。平面和空间中坐标变换公式的线性特征——费马不可能不曾注意到它——例如被欧拉（{[}108 a{]}，（1），第 IX 卷，第 II-III 章及附录第 IV 章）加以突出，他在此基础之上建立了平面曲线以及曲面按次数的分类（正是因为这些公式的线性性，次数才是不变的）；也是他（出处同上，第 XVIII 章）引入了“仿射性”一词，用以指称可以经由变换 \(x' = ax, y' = by\) 互相推演出的那些曲线之间的关系（但在这一定义中，他没有察觉到任何几何不变量，因为该定义仍然依附于坐标轴的一个特定选取）。稍后，我们看到拉格朗日（{[}191{]}，第 III 卷，第~661-692~页）以一整部论文——它将长期理所当然地享有盛名——专门研究三维解析几何中典型的线性和多重线性问题。大约在同一时期，联系到求过若干给定点的平面曲线这一线性问题，行列式的概念成形了，最初以某种多少是经验的方式，见于克拉默 {[}72{]} 和贝祖 {[}21{]}；这一概念随后被多位作者加以发展，其本身及其本质性质最终由柯西（{[}56 a{]}，（2），第 I 卷，第~91-169~页）和雅可比（{[}171{]}，第 III 卷，第~355-392~页）明确地提炼了出来。

另一方面，当数学家们对一次方程还多少有些轻视之际，微分方程的求解却相反构成一个重大问题；在这些方程之中，人们自然很早就挑出了线性方程——不论其系数是否为常数——对它们的研究有助于提升线性性及与之相伴诸物的价值。达朗贝尔 {[}75 c{]}、拉格朗日（{[}191{]}，第 I 卷，第~471-668~页）和欧拉（{[}108 a{]}，（1），第 XII 卷）的情形正是如此；但其中只有第一位认为值得说出：非齐次方程的通解是一个特解与相应的齐次方程的通解之和；此外，当这些作者断言 n 阶线性齐次方程的通解是 n 个特解的线性组合时，他们并没有补充说这些特解必须是线性无关的，也不曾费力把后一概念弄明确；在这一点上，正如在许多其他问题上一样，看来只有柯西在综合理工学校的教学才澄清了问题（{[}56 b{]}，第~573-574~页）。但是拉格朗日（出处同上）已经同样引入了（诚然仅仅是为了计算，而且没有给它命名）伴随方程 \(L^{*}(y) = 0\)，即线性微分方程 \(L(y) = 0\) 的伴随，这是对偶性的一个典型例子，其依据是关系式

\[
\int z L (y) d x = \int L ^ {*} (z) y d x
\]

它在 y 与 z 于积分区间两端为零时成立；更确切地说，在高斯明确给出 3 个变量的线性代换的转置之前 30 年，我们在这里无疑看到了“函数算子”的第一个例子：\(L^{*}\) 是借助一个双线性函数（这里是积分 \(\int yzdx\)）给出的算子 L 的转置或“伴随”。

与此同时，同样是在拉格朗日那里（{[}191{]}，第 III 卷，第~695-795~页），线性代换——起初是 2 个或 3 个变量的——开始向算术进军。显然，函数 \(F(x, y)\) 当 \(x\) 与 \(y\) 取遍全部整数值时所取的值组成的集合，在对 \(x, y\) 施行任意一个具整系数且行列式为 1 的线性代换时并不改变；拉格朗日正是以这一基本观察为基础，建立了数用型表示的理论以及型的化简理论；而高斯则以一步——我们如今已难以完全估量其全部大胆之处——由此提炼出等价的概念和型的类的概念（参见第~50~页）；在这一点上，他认识到需要若干关于线性代换的初等原理，并特别引入了转置或伴随的概念（{[}124 a{]}，第 I 卷，第~304~页）。从那时起，2 个、3 个以及后来的 \(n\) 个变量的二次型的算术研究和代数研究，与它们紧密相连的双线性型的研究，以及更近时期这些概念向无穷多个变量的推广，直到我们这个时代，一直是线性代数最富有成果的进步源泉之一（参见第~136~页以下）。

但是，或许更为决定性的一步是，高斯在这部《算术研究》中创立了（参见第~51~页）有限阿贝尔群的理论，它以四种不同的方式出现：整数关于模 m（m 为整数）的加法群的方式，与 m 互素的数关于模 m 的乘法群的方式，二元二次型的类群的方式，以及最后，m 次单位根的乘法群的方式；而且，如我们已经指出的，高斯显然是把所有这些群当作阿贝尔群、或者更确切地说当作 Z 上的模来处理的，研究它们的结构、它们的同构关系等等。在“复整数”\(a + bi\) 的模中，我们后来看到他研究一个 Z 上的无限模，他无疑没有忽略它与椭圆函数的周期模（由他在复域中发现）的同构；无论如何，这一思想在雅可比那里已经清楚地出现，例如在他关于三个周期的函数不可能性的著名证明中，以及在他关于阿贝尔积分反演问题的见解中（{[}171{]}，第 II 卷，第~25-50~页），并很快以克罗内克的诸定理告终（{[}186 a{]}，第 III\(_{1}\) 卷，第~49-109~页）。

在这里，在我们一直试图追寻其脉络、有时还有其迂回的这几股潮流之中，又有另一股长期潜行于地下的潮流汇入。正如别处将更详细地表明的那样（见第~130~页），前一世纪所理解意义下的“纯粹”几何学，也就是说，本质上是不用坐标的平面和空间的射影几何，是 17 世纪由德萨格 {[}84{]} 创立的，他的思想曾被一位费马赏识其真正价值、被一位帕斯卡付诸使用，随后便归于废弃，被解析几何的辉煌进步所遮蔽；到 18 世纪末，它由蒙日，随后由彭赛列及其追随者布里昂雄、沙勒，重新带回荣耀的地位，有时完全地、有意地与解析方法相分离，有时（特别是在德国）又与后者紧密地混合在一起。然而，射影变换，无论从哪种观点（综合的或解析的）加以考虑，当然无非是射影坐标或“重心”坐标上的线性代换；圆锥曲线（在 17 世纪）、以及后来的二次曲面——它们的射影理论长期是这一学派的主要研究对象——无非就是二次型，而我们上文已经揭示过二次型与线性代数的密切联系。与这些概念相汇合的还有配极的概念：极点与极线的理论同样由德萨格创立，它在蒙日及其后继者手中、不久又以对偶原理之名，成为变换几何定理的一种有力工具；即使人们不敢断言它与伴随微分方程的联系只是很晚才被注意到（平谢莱在世纪末指出了这种联系），有一点毕竟没有被错过——沙勒对此作证（{[}60 b{]}，第~55~页）——那就是注意到它与球面三角形对偶（互反球面三角形）概念的亲缘关系，后者早在 16 世纪就由韦达（{[}319{]}，第~418~页）和斯涅耳引入球面三角学。但是射影几何中的对偶性只是向量空间对偶性的一个方面，这时要考虑到从仿射空间过渡到射影空间（它是仿射空间对“数乘”关系的商空间）所强加的那些修正。

19 世纪比我们历史上的任何时期都更富有一流的数学家；而要在寥寥几页之中——即使只限于最突出的方面——描述这些思想运动在他们手中的融合所产生的一切，是很困难的。一方面是纯粹综合的方法，它像一个普罗克鲁斯提斯之床，其正统信徒们在其中自我折磨，另一方面是与任意强加于空间的坐标系相联系的解析方法，人们很快就感到需要一种几何演算，它被莱布尼茨梦想过而未创立，被卡诺 {[}51{]} 不完善地勾勒过：首先出现的是向量的加法，它隐含于高斯对虚数的几何表示及他由此对初等几何所作的应用之中（参见第~162~页），由贝拉维蒂斯以“等偏差法”之名加以发展，并在格拉斯曼、默比乌斯、哈密顿那里取得其确定的形式；与此同时，默比乌斯以“重心演算”之名给出了它的一个适于射影几何之需的版本（{[}223{]}，第 I 卷）。

与此同时，在同样一些人中间，实现了从平面和“通常”空间到 n 维空间的过渡，这一过渡是如此自然（一旦走上这条道路），以至于我们已经看到它被费马预告过；这一过渡甚至是不可避免的，因为对于 2 个或 3 个变量可用几何术语自身来解释的代数现象，对任意多个变量都毫无改变地成立；因此，若在使用几何语言时强加一个限于 2 维或 3 维的限制，对现代数学家来说将是一个桎梏，其碍事程度不亚于那个一直阻止希腊人把数的概念扩张到不可公度之量的比的桎梏。也正因为如此，与 n 维空间有关的语言和观念几乎同时在各方面出现，在高斯那里是隐约的，在下一代数学家那里则是清晰的；而他们在使用这些观念时或大或小的大胆程度，与其说取决于他们的数学倾向，不如说取决于他们的哲学见解、甚至仅仅实际的考虑。无论如何，凯莱和格拉斯曼在大约 1846 年就以极大的自如驾驭这些概念（凯莱在与格拉斯曼相异之处说道（{[}58{]}，第 I 卷，第~321~页），“无需诉诸任何形而上学的概念”）；在凯莱那里，人们始终非常接近于解析解释和坐标，而在格拉斯曼那里，从一开始——随着 n 维空间中向量的加法——领先的是几何方面，并最终引出我们稍后将要谈到的那些发展。

然而，来自高斯的推动正以两种不同的方式把数学家引向代数即超复系统的研究。一方面，人们不可避免地要尝试以引入“虚数单位”\(i = \sqrt{-1}\) 之外的其他方式来扩张实数的范围，从而也许能开辟出比复数域更广阔、更富成果的领域。高斯使自己确信（{[}124 a{]}，第 II 卷，第~178~页）这种扩张是不可能的，至少当人们想保留复数的主要性质、也就是说用现代的语言讲、使它成为一个交换域的那些性质时是如此；而他的同时代人，或由于他的影响，或独立地，似乎也分享了这一信念，它直到很久以后才由魏尔斯特拉斯（{[}329 a{]}，第 II 卷，第~311-332~页）以一个精确的定理加以证明。但是，一旦把复数的乘法解释为平面中的旋转，人们若想把这一想法推广到空间，就会被引导去考虑（由于空间中的旋转构成一个非交换群）非交换的乘法；这正是引导哈密顿 \(^{1}\) 发现四元数 {[}145 a{]} 的思想之一，四元数是除环的第一个例子。这一例子的独特性（正如弗罗贝尼乌斯后来所证明的，它是在实数域上可以构造出来的唯一例子）在某种程度上限制了它的影响范围，尽管有、或者甚至恰恰由于形成了一个狂热的“四元数主义者”学派：一种奇怪的现象，它后来在格拉斯曼的著作周围重演，随后又有那些普及者从哈密顿和格拉斯曼那里提炼出所谓“向量演算”。稍后格雷夫斯和凯莱放弃结合律而构造出“凯莱数”，也并未开辟出一条有意义的矿脉。但是，在西尔维斯特引入了矩阵并（未给它命名）明确定义了秩（{[}304{]}，第 I 卷，第~145-151~页）之后，又是凯莱（{[}58{]}，第 II 卷，第~475-496~页）创立了矩阵的演算，而且他不曾忽略（这是一个本质的事实，后来常被忽视）指出：矩阵无非是线性代换的一种缩写记号，正如事实上高斯用 \((a, b, c)\) 表示型 \(aX^2 + 2bXY + cY^2\) 一样。这里的情形除了是西尔维斯特和凯莱关于行列式及与之相关的一切的大量论著的一个方面——对我们来说无疑是最有趣的一个方面——之外，别无其他，那些论著布满了巧妙的恒等式和令人印象深刻的计算。

格拉斯曼发现的（除其他东西外）还有实数域上的一个代数，即与他的名字联系在一起的外代数。他的著作（{[}134{]}，第 I₁ 卷）甚至先于哈密顿的著作，是在一种几乎完全彻底的精神孤独中创造出来的，长期鲜为人知，这无疑是由于它的独创性，也由于它开头笼罩着的哲学迷雾——例如这迷雾一开始就使默比乌斯望而却步。格拉斯曼怀着与哈密顿类似但视界更为宽广的关切（而且他很快就意识到，这些关切正是莱布尼茨本人的关切），建造起一座庞大的代数—几何大厦，它奠基于 n 维向量空间的一种几何的或者说“内蕴”的观念（差不多已经公理化）之上；在他最终得到的结果之中，最初等的一些包括：向量线性无关的定义、维数的定义以及基本关系式 dim V + dim W = dim (V + W) + dim (V ∩ W)（出处同上，第~209~页；参见{[}134{]}，第 I₂ 卷，第~21~页）。但首先是多向量的外乘法、继之内乘法，为他提供了工具，使他得以从容地处理线性代数本身的问题，然后是涉及欧氏结构即向量正交性的那些问题（在后者中他找到了他所缺少的对偶性的等价物）。

高斯在超复系统研究中开辟的另一条道路，是从复整数 \(a + bi\) 出发的那一条；从这些整数出发，人们很自然地继续走向整数环 Z 上或有理数域 Q 上的更一般的代数或超复系统，首先走向高斯已经预见到的、由单位根生成的那些系统，然后走向代数数域和代数整数模：这些正是库默尔著作的主要对象，而对前者的研究则由狄利克雷、埃尔米特、克罗内克、戴德金着手进行。在这里，与实数域上的代数的情形相反，无须放弃交换域的任何一条特征性质，而整个 19 世纪的工作也正是限于这些系统。但那些线性的性质，例如寻求域中整数的基（它对于判别式的一般定义是不可或缺的），在许多方面都扮演着本质的角色；而且，无论如何在戴德金那里，那些方法注定要变成典型“超复的”；此外，戴德金本人虽然并未以一般的方式处理代数的问题，却意识到他工作的这一方面，意识到例如把它与魏尔斯特拉斯关于实数域上超复系统的结果联系起来的东西（{[}79{]}，特别是第 II 卷，第~1-19~页）。与此同时，代数数域中单位乘法群结构的确定——由狄利克雷在其著名的论文（{[}92{]}，第 I 卷，第~619-644~页）以及几乎同时由埃尔米特（{[}159{]}，第 I 卷，第~159-160~页）完成——极适合于澄清关于 Z 上的模、它们的生成系以及（当它们具有时）它们的基的观念。随后是理想的概念：戴德金在代数数域中把它定义为（域的整数环上的）一个模（{[}79{]}，第 III 卷，第~251~页），而克罗内克则在多项式环中（以“模系”之名）引入一个等价的概念（{[}186 a{]}，第 II 卷，第~334-342~页），这给出了 Z 以外更一般的环上的模的第一批例子；而且，在这些相同的作者那里、然后在希尔伯特那里，带算子群的概念以及在特殊情形下从这样一个群出发总能构造出一个适当定义的环上的模的可能性，被一点一点地释放出来。

与此同时，对二次型和双线性型的算术—代数研究及其“化简”（或者换一种说法，矩阵及其“不变量”的化简），导致发现了关于线性方程组求解的一般原理，这些原理曾因缺少秩的概念而未被雅可比注意到。\(^{2}\) 整系数线性方程组的整数解问题被提出并解决，起初是由埃尔米特在一个特殊情形中，然后由 H. J. 史密斯以其全部一般性解决（{[}287{]}，第 I 卷，第~367-409~页）；后者的结果直到 1878 年才由弗罗贝尼乌斯重新发现，那是在克罗内克所制定、魏尔斯特拉斯也参与其中的一个庞大的研究纲领的框架之内；顺便说一句，正是克罗内克在研究过程中给出了实（或复）系数线性方程组诸定理的最终形式，这些定理也由《爱丽丝漫游奇境记》的著名作者在一部不起眼的教材中以他所特有的缜密细心加以阐明；至于克罗内克，他不屑于发表自己的结果，而把它留给同事和弟子们；甚至连“秩”这个词也只是由弗罗贝尼乌斯引入的。同样是在柏林大学的教学中，克罗内克 {[}186 b{]} 和魏尔斯特拉斯给出了行列式的“公理化”定义（即 n 维空间中 n 个向量的一个多线性反对称函数，规范化为在单位矩阵上取值 1），这一定义等价于可以从格拉斯曼的计算中推导出来的那个定义；也还是在他的讲义中，克罗内克在仍未感到需要给它命名、且尚非内蕴形式的情况下，引入了空间的张量积和矩阵的“克罗内克”积（由施于各因子的给定线性代换在张量积上诱导出的线性代换）。

这些研究也不应与不变量理论相分离，后者由凯莱、埃尔米特和西尔维斯特创立（埃尔米特后来在书信中称之为“不变量三位一体”），而从现代的观点看，它首先是一种线性群的表示理论。作为射影几何对偶性的代数等价物，出现了同步变量列与逆步变量列的区分，也就是说，一个空间中的向量与对偶空间中的向量的区分；在首先注意到 2 个或 3 个变量的低次型、然后是任意次数的型之后，人们很快便考察起多组“同步”或“逆步”变量的双线性型、继而多线性型来，这等价于张量的引入；当里奇和列维—奇维塔在不变量理论的启发下于 1900 年把“张量分析”引入微分几何 {[}257{]} 时，这一点便变成了自觉的并得到普及——张量分析后来因“相对论者”物理学家们的使用而风靡一时。也还是不变量理论、微分几何以及偏微分方程理论（特别是所谓普法夫问题及其推广）的逐步相互渗透，逐渐引导几何学家们首先考虑微分的双线性反对称型，特别是一次型的“双线性协变量”（1870 年由李普希茨引入，随后由弗罗贝尼乌斯研究），最终导致 E. 嘉当（{[}52 a{]}，第 II₁ 卷，第~303-396~页）和庞加莱（{[}251 b{]}，第 III 卷，第 XXII 章）创立外微分形式的演算。庞加莱是在构成其积分不变量的脉络中，把它们作为出现在多重积分中的表达式引入的，而嘉当——无疑是在他关于代数的研究的引导之下——则以一种不那么形式的方式引入它们，但也不曾放过指出其演算的代数部分与格拉斯曼的外乘法完全相同（他采用的名字即由此而来），从而最终把后者的著作放回到它真正的位置上。把外微分形式翻译成张量分析的记号，此外还立即表明了它们与反对称张量的联系，而一旦回到纯代数的观点，这就会使人看出：它们之于交错多线性型，正如协变张量之于任意多线性型；这一方面随着现代的线性群表示理论而变得更加清晰；由此人们便认识到，例如，魏尔斯特拉斯和克罗内克给出的行列式定义与由格拉斯曼的演算得出的定义实质上相同。

我们于是到达了现代，在那里，公理化方法和结构概念（起初被感受到，直到最近才被定义）使我们得以把直到那时还不可分解地混杂在一起的概念分离开来，把模糊的或潜意识的东西表述出来，以适当的普遍性去证明那些此前只在特殊情形中已知的定理。皮亚诺是公理化方法的创立者之一，也是最早按其真正价值赏识格拉斯曼著作的数学家之一，他早在 1888 年（{[}246 b{]}，第 IX 章）就给出了实数域上向量空间（不论是否有限维）的公理化定义，并以一种完全现代的记号给出了从这样一个空间到另一个空间的线性映射；稍后，平谢莱试图把这样构想出来的线性代数应用于函数论，其所取的方向诚然几乎没有显示出多产的迹象；至少他的观点使他得以在“拉格朗日伴随”中认出线性映射的转置的一个特殊情形：这一点不久便更清晰地出现，而且既出现在微分方程中，也出现在偏微分方程中，那是在希尔伯特及其学派关于希尔伯特空间及其分析应用的令人难忘的工作过程中。正是联系着后面这些研究，托普利茨 {[}309 a{]} 在（借助于坐标）引入实数域上最一般的向量空间的同时，作出了如下基本观察：行列式理论对于线性代数主要定理的证明是没有用处的，这就使得这些定理可以毫无困难地推广到无限维空间；他还指出，这样理解的线性代数当然适用于任意的交换基域。

另一方面，随着巴拿赫在 1922 年引入以他的名字命名的空间，\(^{3}\) 人们遇到了——诚然是在一个既属拓扑又属代数的问题中——与其对偶不同构的空间（见第~214~页以下）。在一个向量空间与其对偶空间之间，本来就不存在“典范”同构，也就是说由结构决定的同构，这一点长期以来就反映在同步与逆步的区分之中。尽管如此，似乎毫无疑问的是，一个空间与其对偶之间的区分，只是随着巴拿赫及其学派的工作才最终确立起来；也正是这一工作唤起了对余维数概念的兴趣。至于一个空间的向量子空间与其对偶空间的向量子空间之间的对偶性或“正交性”，它今天的表述方式显示出与伽罗瓦理论基本定理的现代表述、与局部紧阿贝尔群中所谓庞特里亚金对偶性的一种绝非偶然的类似；这可以追溯到韦伯，他在算术研究的过程中于 1866 年就有限群陈述了其基础 {[}327 d{]}；在伽罗瓦理论中，子群与子域之间的“对偶性”随戴德金和希尔伯特而成形；而向量子空间之间的正交性，则明显地一方面派生于射影几何中线性流形之间的对偶性，另一方面派生于欧氏空间或希尔伯特空间中完全正交流形的概念及其性质（它的名字即由此而来）。所有这些线索在当代汇合到代数学家如 E. 诺特、阿廷和哈塞、拓扑学家如庞特里亚金和惠特尼的手中，而且这两组人之间并非没有相互施加影响。

与此同时，人们着手进行一种批判性的审查，其意图是在每一点上消除那些并非真正不可缺少的假设，尤其是那些会妨碍某些应用的假设。由此便实现了在向量空间概念中以环代替域的可能性，而且，通过创立模的一般概念，可以同时处理这些空间、阿贝尔群、已被克罗内克、魏尔斯特拉斯、戴德金、施泰尼茨研究过的那些特殊的模，乃至带算子群，例如把若尔当—赫尔德定理应用于它们；与此同时，借助右模与左模的区分实现了向非交换性的过渡，而美国学派（韦德伯恩、迪克森）特别是德国学派（E. 诺特、阿廷）的代数理论的现代发展正导向这一区分。

最后，相当晚近才出现了我们必须在这里谈及的诸倾向中的最后一种：它以胚胎的形式包含于戴德金的定理（{[}79{]}，第 III 卷，第~29~页）——即交换域的任意多个自同构总是线性无关的——之中，伽罗瓦理论的线性化由阿廷 {[}7 a{]} 完成，随后很快被推广到交换域的任意扩张，甚至非交换域的扩张（{[}172 d{]}，第 VII 章）。

\chapter{多项式与交换域。}

交换域的理论——以及与它紧密相联系的多项式理论——直接来源于直到 19 世纪中叶一直构成经典代数之主要对象的东西：代数方程的求解，以及与之等价的几何作图问题。

一旦着手求解次数 \textgreater{} 1 的代数方程，人们就会遇到一类全新的计算困难：仅凭从给定量出发的“有理”计算，是不可能完成对一个未知数的确定的。这一困难必定在非常早的阶段就被察觉到了；而巴比伦人知道把二次方程和双二次方程的求解化归为一种新的代数运算——开平方——则必须算作他们最重要的贡献之一（流传到我们手中的文本里大量以这种方式求解的数值方程可以证明这一点（{[}232{]}，第~183-193~页））。就形式计算而言，古代在代数方程求解问题上从未超过这一点；古典时代的希腊人确实仅限于用几何的语言重新找到巴比伦人的公式，而以代数形式使用这些公式，在海伦（公元 100 年）和丢番图之前尚无明证。

希腊人是在一个完全不同的方向上取得决定性进展的。关于巴比伦人如何设想和计算非平方整数的平方根，我们所知甚少；\(^{1}\) 在流传到我们手中的涉及这一问题的罕见文本里，他们似乎满足于相当粗糙的近似方法（{[}232{]}，第~33-38~页）。毕达哥拉斯学派曾经严格地确定了可公度量的概念，并赋予它一种近乎宗教的性质，故而不可能持这种观点；而把他们最终引向证明这个数是无理数的，可能正是把 \(\sqrt{2}\) 表为有理数的反复尝试的失败。\(^{2}\)

我们将在别处（参见第~148~页）叙述这一标志着数学史重大转折的发现如何深刻地反作用于希腊人的“数”的概念，并如何引导他们创立一种纯几何特征的代数，以便在其中找到一种表示不可公度之比（或者也许是为这种比给出一种“存在性”证明）的办法——他们拒绝把这些比当作数。在大多数情形，他们把一个代数问题化归为两条适当选取的辅助平面曲线的相交，或者化归为这类相交的若干次相继确定。一些迟出的、权威性甚小的传统把这类作图的一种最初的分类——它注定有长久而辉煌的历程——归于柏拉图：看来是出于比数学更偏哲学的理由，他把所谓“直尺和圆规”的作图，也就是说除直线和圆之外不用任何辅助曲线的那些作图，单独放到一边。\(^{3}\) 无论如何，欧几里得在其《原本》{[}107{]} 中仅限于专门处理可用这种方式求解的问题（然而并未以一个特别的名称来刻画它们）；这一情况无疑大大有助于把此后几个世纪数学家的注意力固定在这些问题上。但我们今天已经知道，\(^{4}\) 可以用“直尺和圆规”求解的代数方程属于一种非常特殊的类型；特别地，三次不可约方程（在有理数域上）不能以这种方式求解，而希腊人很早便遇到了一些一直享有盛名的三次问题，例如倍立方（方程 \(x^{3}=2\) 的解）和三等分角；化圆为方则相反使他们面对一个超越的问题。我们看到，为了解决这些问题，他们引入了许多代数的（圆锥曲线、狄奥克莱斯的蔓叶线、尼科墨德斯的蚌线）或超越的（狄诺斯特拉图的割圆曲线、阿基米德螺线）曲线；这不可能不把他们引向对这些不同曲线的自主研究，从而为解析几何、代数几何和无穷小演算未来的发展准备了道路。但这些方法对代数方程的求解没有带来任何进展，\(^{5}\) 而古代唯一对这一问题作出了像样贡献、并对中世纪和文艺复兴时期的代数学家施加了持久影响的著作，就是欧几里得《原本》{[}107{]} 的第 X 卷；在这一卷中（有些史家想把其中的主要结果归于泰阿泰德），他考

察由若干根式的组合所得的表达式，例如 \(\sqrt{\sqrt{a}} \pm \sqrt{b}\)（a、b 为有理数），给出这些表达式为无理数的条件，把它们分成许多类别（他证明这些类别互不相同），并研究这些各种各样的无理数之间的代数关系，例如我们今天会写成

\[
\sqrt {\sqrt {p} + \sqrt {q}} = \sqrt {1 / 2 (\sqrt {p} + \sqrt {p - q})} + \sqrt {1 / 2 (\sqrt {p} - \sqrt {p - q})};
\]

这一切都是用《原本》惯用的几何语言表述的，这使得阐述格外沉闷而笨拙。

在古典希腊数学衰落之后，有关代数方程的各种概念发生了变化。毫无疑问，在整个古典时期，希腊人掌握着平方根的不定近似方法，可惜我们对之了解得很差。\(^{6}\) 随着印度人，然后是中世纪的阿拉伯人及其西方效仿者，各次开方变成了一种倾向于被视为与代数的有理运算同等基本的运算，并且像后者一样，用越来越便于计算的记号来表示。\(^{7}\) 在整个中世纪中日臻完善的二次方程理论（根的个数、负根、不可能的情形、二重根）以及双二次方程理论，给出了方程“根式解”公式的范本，几个世纪里代数学家们都试图从中仿造出解更高次方程、首先是三次方程的类似公式。13 世纪把阿拉伯科学介绍到西方的主要人物比萨的莱昂纳多认识到，无论如何，欧几里得在其第 X 卷中分类的那些无理数不能为此目的所用（在这样一门包含如此多不可能性的理论中，这又是一个不可能性的证明），而且我们看到他已经在尝试用立方根做类似的计算，得到了诸如

\[
\sqrt [ 3 ]{16} + \sqrt [ 3 ]{54} = \sqrt [ 3 ]{250},
\]

这样的关系式，它们与欧几里得的平方根公式类似（此外，更早的例子在阿拉伯人那里已经可以找到）。但还要再经过三个世纪徒劳的努力，费罗才在 16 世纪初最终对方程 \(x^{3} + ax = b\) 得到了求解公式 \(^{8}\)

\[
x = \sqrt [ 3 ]{\frac {b}{2} + \sqrt {\left(\frac {b}{2}\right) ^ {2} + \left(\frac {a}{3}\right) ^ {3}}} + \sqrt [ 3 ]{\frac {b}{2} - \sqrt {\left(\frac {b}{2}\right) ^ {2} + \left(\frac {a}{3}\right) ^ {3}}}.\tag{5.1}
\]

我们在这里无法描述这一轰动性发现之富有色彩的一面——它一方面在塔尔塔利亚与另一方面在卡尔丹及其学派之间激起的争吵——也无法描述作为主角的学者们往往颇为动人的形象。但我们必须指出，随后在卡尔丹及其学生手中，方程理论取得了决定性的进步。卡尔丹对使用负数的反感比他大多数同时代人要少，他由此观察到三次方程可以有三个根，而双二次方程有四个根（{[}50{]}，第 IV 卷，第~259~页），并且他指出，\(x^3 + bx = ax^2 + c\)（他知道如何消去这个方程中的 \(x^2\) 项）的三个根之和总等于 \(a\)（出处同上）。无疑是在这一关系以及对其普遍性质的直觉的引导下，他萌生了根的重数这一概念的最初想法；尤其是，他壮起胆子（并非没有一番言辞上的铺垫）对含负数平方根的表达式进行形式计算。他很可能是被如下事实引到这一步的：当 \(\left(\frac{b}{2}\right)^{2}+\left(\frac{a}{3}\right)^{3}<0\) 时（这是一种被称为“不可约”的情形，卡尔丹已经认识到其中存在三个实根），这样的表达式在公式（1）的使用中会自然出现；无论如何，这一点在他的学生 R. 邦贝利那里清楚地显现出来，邦贝利在其《代数学》（{[}28 a{]}，第~293~页）中证明了关系式

\[
\sqrt [ 3 ]{2 + \sqrt {-121}} = 2 + \sqrt {-1}
\]

并且注意以与现代的阐述十分接近的形式明确给出复数的计算规则。\(^{9}\) 最后，早在 1545 年，卡尔丹的另一个学生 L. 费拉里就借助一个三次的辅助方程成功地解出了普遍的四次方程。\(^{10}\)

在如此迅速的进展之后，直到 18 世纪中叶的随后一个时期，几乎只是发展意大利学派引入的那些新观念。得益于韦达为代数记号带来的这一本质进步，他得以用一般的方式表达代数方程的系数与根之间的关系，至少当诸根皆为正数时是如此 \(^{11}\)（{[}319{]}，第~158~页）。更大胆的 A. 吉拉尔 {[}129{]} 则毫不犹豫地断言（当然没有证明）：只要把“不可能的根”各按其重数计入，n 次方程就恰好有 n 个根，而且这些根满足韦达给出的那些关系；他还第一次得到了根的幂和直到指数 4 的表达式。

但 17 世纪的精神转向了其他方向，代数只是作为间接的后果，才从解析几何和无穷小演算的新发现中得到些许好处。于是，与笛卡儿求代数曲线切线的方法（参见第~178~页）相联系，出现了代数方程之根的重数判别法，它由他的学生许德（{[}85 b{]}，第 I 卷，第~433~页及第~507-509~页）陈述。代数函数与超越函数之间的区分，无疑也须追溯到笛卡儿的影响，它平行于他在其《几何学》中就“几何”曲线与“机械”曲线所作的区分（参见第~176~页和第~193~页）。无论如何，这一区分在 J. 格雷戈里那里是完全清楚的，他在 1667 年甚至试图证明圆扇形的面积不可能是由弦和半径表出的代数函数（{[}136 a{]}；参见{[}155{]}）。“超越的”一词出自莱布尼茨，他在其整个生涯中从未间断对这类分类问题的兴趣，并在大约 1682 年发现了格雷戈里所寻求结果的一个简单证明，即证明了 sin x 不可能是 x 的代数函数（{[}198 a{]}，第 V 卷，第~97-98~页）。\(^{12}\) 他与他的朋友奇恩豪斯，是他那个时代仅有的几位仍对代数方程“根式解”问题感兴趣的数学家。在其早期阶段，我们看到他研究三次方程的“不可约情形”，并确信（事实上并无充分的证明）在这种情形下不可能从求解公式中除去虚量（{[}198 d{]}，第~547-564~页）。大约在同一时期，他也未能成功地攻克五次方程的根式解问题；后来，当奇恩豪斯声称借助形如 \(y = P(x)\) 的变换（其中 P 是一个适当选取的四次多项式）使方程中除两端两项外的所有项都消失、从而解决了问题时，莱布尼茨立即看出，确定 \(P(x)\) 的系数的那些方程的次数 \textgreater{} 5，并断定这个方法注定要失败（{[}198 d{]}，第~402-403~页）。

看来，正是新分析学的需要使人们对代数的兴趣逐渐恢复起来。由莱布尼茨和约翰·伯努利完成的有理函数的积分法，以及与之紧密相联系的虚对数问题，成为深化虚数运算的契机，也使多项式分解为一次因子的这一问题（“代数基本定理”）得到回答。\(^{13}\) 从 18 世纪初起，二项方程 \(x^{n}-1=0\) 的求解就被科茨和棣莫弗化归为把圆分成 n 等份；因此，要得到其根的“根式”表达式，只需会做 n 为奇素数的情形，而棣莫弗注意到，代换 \(y = x + \frac{1}{x}\) 那时便把问题化为解一个 \((n - 1)/2\) 次方程的“根式解”。至于“基本定理”，在一般解法的“根式解”屡次受挫（包括欧拉的若干次尝试（{[}108 a{]}，（1），第 VI 卷，第~1-19~页和第~170-196~页））之后，人们开始寻求先验的、不使用显式求解公式的证明。不必深入所提诸方法的细节（它们最终引出拉格朗日和高斯的证明；参见第~91~页和第~161~页），这里应当指出 18 世纪中叶人们看待这一问题所取的观点：人们假定（除了对一般性的一种模糊感觉——它无疑与 A. 吉拉尔那里一样，来自系数与根之间存在关系这一事实——之外，没有任何别的理由）n 次方程总有 n 个“理想”根，可以对它们像对数那样进行计算，却不知道它们是否为（实的或复的）数；而必须证明的则是（必要时可使用关于理想根的计算）：这些根中至少有一个是通常的复数。\(^{14}\) 以这种有缺陷的形式，人们认识到这里已经有了“形式添加”这一一般观念的第一个胚胎，它不顾高斯的反对（{[}124 a{]}，第 III 卷，第~1~页），将成为现代交换域理论的基础。

随着拉格朗日（{[}191{]}，第 III 卷，第~205-421~页）和范德蒙德 {[}315{]} 的那些基本论文，1770 年开启了代数方程理论历史上一个新的决定性时期。继直到那时一直无可争议地占统治地位的、多少有些碰运气的求解公式尝试的经验主义之后，随之而来的是对所述问题以及可望解决它们的方法的系统分析，这一分析在六十年后将引出伽罗瓦的最终结果。拉格朗日和范德蒙德都从根式的多个取值在次数 \(\leq 4\) 的方程求解公式中引入的多义性出发；这一事实曾引起欧拉的注意（{[}108 a{]}，（1），第 VI 卷，第~1-19~页），他特别指出了在费罗的公式中应当如何把其中出现的各根式的取值适当配对，以便得到 3 个根而不是 9 个。\(^{15}\) 拉格朗日注意到，费罗公式中的每个立方根都可以写成 \(\frac{1}{3}(x_1 + \omega x_2 + \omega^2 x_3)\) 的形式，其中 \(\omega\) 是 1 的一个立方根，\(x_1, x_2, x_3\) 是所论方程的三个根，按某种次序取出；他作出了基本的观察：三个根的函数 \((x_1 + \omega x_2 + \omega^2 x_3)\) 在三个根的所有置换下只能取两个不同的值，这就先验地解释了求解这个方程的那些方法何以成功。对四次方程求解方法的类似分析把他引到四个根的函数 \(x_{1}x_{2} + x_{3}x_{4}\)，它在根的所有置换下只取三个不同的值，因而是以所给方程的系数的有理函数为系数的一个三次方程的根；\(^{16}\) 拉格朗日说，这些事实构成了“真正的原则，以及所谓三次和四次方程求解的形而上学 \(^{17}\)”（{[}191{]}，第 III 卷，第~357~页）。凭借这些例子，他提议一般地对 n 次方程研究有理函数 V 在 n 个根被任意置换时所能取的值的个数 \(\nu\)；\(^{18}\) 实际上，他由此开创了（用的是仍紧密依附于方程论的术语）群的理论和域的理论，并且已经用与今天所用的相同的原则得到了它们的几个基本结果。例如，他借助今天用来证明有限群的子群的阶整除群的阶的那一论证，证明了数 \(\nu\) 是 n! 的一个因子。更为引人注目的是这样一个定理：若 \(V_{1}\) 和 \(V_{2}\) 是根的两个有理函数，且 \(V_{1}\) 与 \(V_{2}\) 在同样的置换下不变，则其中每一个都是另一个与方程的系数的有理函数（这是伽罗瓦把伽罗瓦扩张的子扩张刻画为其伽罗瓦群的不变量域的定理的一个特殊情形）：他说，“这个问题在我看来是方程论中最重要的问题之一，我们将给出的它的一般解，将是使代数的这一部分焕然一新的手段”（{[}191{]}，第 III 卷，第~374~页）。

所有这些研究在拉格朗日的心目中自然只是准备步骤，为的是分析通过逐次化归为次数较低的方程来求解代数方程的可能方法，而他证明了这样的方法联系于构造在根的置换下取值少于 n 个的根的有理函数。无疑是在他关于三次方程的结果的引导下，他一般地引入了“拉格朗日预解式”\(y_{k} = \sum_{h=1}^{n} \omega_{k}^{h} x_{h}\)，其中 \(\omega_{k}\) 是一个 n 次单位根（\(1 \leq k \leq n\)），他清楚地表明这 n 个数的知识如何导出诸根 \(x_{k}\) 的知识，并一般地寻求 \(y_{k}\) 所满足的方程的次数；例如他证明，若 n 为素数，则诸 \(y_{k}\) 是一个 n-1 次方程的根，其系数是某个 \((n-2)!\) 次方程的一个根的有理函数，而该方程的系数又可由所给方程的系数有理地表出。他总结道：“如果我没有弄错的话，这些就是方程求解的真正原则，以及达到这一目标的最正确的分析；正如可以看到的，一切都可以归结为一种组合计算，借助它，人们可以先验地得到所应期待的结果”（{[}191{]}，第 III 卷，第~403~页）。

至于范德蒙德的论文，它独立于拉格朗日的论文，却在许多点上与后者相遇，特别是在寻求在根的置换下取尽可能少的不同值的根的有理函数这一观念方面，\(^{19}\) 以及他为此目的同样引入的“拉格朗日预解式”的研究方面。他的著作远不具有拉格朗日著作的那种清晰性和普遍性；然而在一点上它明显走得更远，即把同样的观念应用于 n 为奇素数的分圆方程 \(x^{n}-1=0\)。拉格朗日满足于提醒读者，这个方程可化归为一个 \(m=(n-1)/2\) 次的、以有理数为系数的方程，而当 \(n\geq11\) 时并不去求解它；范德蒙德则断言，由于 \(x^{n}-1=0\) 的各个根之间的关系，这个方程的拉格朗日预解式的 m 次幂是有理数；但他仅限于在 n=11 的情形核验这一断言的真实性，而未以一般的方式证明它。

只是到了 30 年后，范德蒙德所预告的结果才由 C. F. 高斯完全证明。\(^{20}\) 他关于方程 \(x^n - 1 = 0\)（\(n\) 为奇素数）的决定性结果，属于他那令人难忘的算术研究的总纲领（{[}124 a{]}，第 I 卷，第~413~页以下），并以一种非常特别的方式显示出他驾驭我们今天所谓的循环群理论的娴熟技艺。在证明了多项式 \(\Phi_n(x) = (x^n - 1)/(x - 1)\) 对奇素数 \(n\) 不可约 \(^{21}\) 之后，他想到把它的 \(n - 1\) 个根写成 \(\zeta^{g^k} = \zeta_k (0 \leq k \leq n - 2)\) 的形式，其中 \(g\) 是同余式 \(z^{n-1} \equiv 1 \pmod{n}\) 的一个原根）（用现代的语言说，这相当于挑明 \(\Gamma\) 即方程 \(\Phi_n(x) = 0\) 的群是循环群这一事实）。对每一个因子 \(e\)（即 \(n-1\) 的因子），他对应地取 \(f = (n-1)/e\) 个“周期”\(\eta_\nu = \xi_\nu + \xi_{\nu+f} + \xi_{\nu+2f} + \cdots + \xi_{\nu+(e-1)f} (1 \leq \nu \leq f)\)，并实质上证明了：诸 \(\eta_{\nu}\) 的以有理数为系数的线性组合构成一个域，它由 \(f\) 个周期 \(\eta_{\nu}\) 中的任何一个生成，并且它在有理数域上的次数为 \(f\)（这个域自然对应于 \(\Gamma\) 的 \(e\) 阶子群）。我们在这里无法深入他的分析的细节以及由此得出的重要算术推论；只需注意到，它特别使他得到了那个关于可以用“直尺和圆规”作出边数等于形如 \(2^{2^k} + 1\) 的素数的多边形的著名定理。至于方程 \(\Phi_n(x) = 0\) 的根式解，它容易由把周期理论应用于一个拉格朗日预解式的 \(f\) 次幂（预解式为 \(\sum_{\nu=0}^{f-1} \omega^\nu \eta_\nu\)，其中 \(\omega^f = 1\)）而得出。

他的同胞鲁菲尼的研究 {[}265{]} 与《算术研究》同时代，而且直接承自拉格朗日；他在拉格朗日停下的地方重新拾起这一问题，提出的目标是证明“一般”五次方程 \(^{24}\) 的“根式解”之不可能性。鲁菲尼的证明冗长而晦涩，尽管几经修改仍不完全；但它已经非常接近阿贝尔后来所得到的（原则上正确的）证明。\(^{25}\) 它的主要意义尤其在于引入了置换的计算以及群论的最初概念，鲁菲尼发展这些内容是为了证明：不存在方程 5 个根的一个函数，当诸根被任意置换时它取的值多于 2 个而少于 5 个。

我们已经说过（见第~51~页），置换群理论的这一最初草稿是如何在几年之后由柯西发展和系统化的。但是，如果说就置换而言、拉格朗日诸观念的发展所必需的那些概念就这样被逐渐澄清了，那么还必须以同样清晰的方式陈述域论的最初原理。这正是鲁菲尼所缺少的东西，而这也正是阿贝尔和伽罗瓦在代数方程求解问题的最后阶段将要做到的。

在他短促的一生中，阿贝尔从未间断过对这一问题的关切。当他还几乎是个孩子的时候，他曾相信自己已经得到了一般五次方程的根式求解公式。后来认识到了自己的错误，他便不达证明出这样的公式并不存在的目的誓不罢休（{[}1{]}，第 I 卷，第~66~页）。但他并未止步于此。当他的竞争者雅可比作为分析学家发展椭圆函数理论时，支配着阿贝尔这方面工作的却是代数的观点，其中心是椭圆函数的分（割）方程理论（{[}1{]}，第 I 卷，第~265、377~页及各处）。他由此用仿照高斯解分圆方程的方法（{[}1{]}，第 I 卷，第~310~页和第~358~页）得到了一些可用根式求解的新类型的方程；\(^{26}\) 由这一结果他进而上升到“阿贝尔”方程的概念，并在一篇著名论文中证明了这类方程可用根式求解（{[}1{]}，第 I 卷，第~478~页）；也正是在这个场合，他精确地定义了在一个给定域（由他所研究的方程的系数生成的域）上不可约的多项式的概念。\(^{27}\) 最后，死亡于 1829 年把他击倒了，那时他正好着手刻划所有可用根式求解的方程的一般问题，并且刚刚向克雷勒和勒让德通报了已经非常接近伽罗瓦的结果（{[}1{]}，第 II 卷，第~219-243、269-270~页和第~279~页）。

三年之后，为这座大厦加冕的任务留给了他 {[}123{]}。像阿贝尔一样，但以一种更为清晰的方式，他一开始就（除术语差异外）定义了一个数属于由给定的诸数所生成的域、添加的概念，以及给定域上的不可约多项式。给定一个方程 \(F(x) = 0\)，它没有重根、系数属于一个给定的域 \(K\)，他相继证明：“总可以构造诸根的一个函数 \(V\)，使得以一切可能的方式在这个函数中置换诸根时所得到的诸值互不相等”，这个函数“具有这样的性质：所提方程的所有根都可以表示为 \(V\) 的有理函数”，以及，设 \(V, V', V'', \ldots\) 是 \(V\) 所满足的不可约方程的诸根，“若 \(a = f(V)\) 是所提方程的一个根，则 \(f(V')\) 同样将是所提方程的一个根”（{[}123{]}，第~47-51~页）；用现代的语言说，他由此证明了 \(V\) 以及它在 \(K\) 上的任一共轭元都生成域 \(N\)，即 \(F\) 的诸根所构成的域。然后他把 \(\Gamma\) 定义为 \(F\) 的群，即诸根 \(x_i\) 的如下置换的集合：把 \(V\)——在各个 \(x_i\) 作为 \(V\) 的函数的有理表达式中——代之以 \(V\) 的任一共轭元，便得到这些置换；他并由此立即得到 \(K\) 中元素凭其在 \(\Gamma\) 的每个置换下不变这一性质的基本刻划（{[}123{]}，第~51~页）。接着他证明，若 \(N\) 包含另一个多项式的根域 \(L\)，则 \(N\) 在 \(L\) 上的群是 \(\Gamma\) 的一个正规子群（这个概念正是他在此时引入的）（{[}123{]}，第~175~页）。最后他由此导出方程可用根式求解的判据，其论证的本质部分如下：设基域 \(K\) 包含所有的单位根，则依假设必存在一个递增序列 \((K_i)_{0 \leq i \leq m}\)，它由介于 \(K\) 与 \(N\) 之间的中间域组成，满足 \(K_0 = K, K_m = N\)，域 \(K_{i+1}\) 是通过向 \(K_i\) 添加一个二项方程 \(x^{n_i} - a_i = 0\)（其中 \(a_i \in K_i\)）的所有根而得到的。因此在 \(\Gamma\) 中应存在一个子群的递减序列 \((\Gamma_i)\)，使得 \(\Gamma_0 = \Gamma, \Gamma_m = \varepsilon\)（中性元），\(\Gamma_{i+1}\) 在 \(\Gamma_i\) 中正规，且商群 \(\Gamma_i / \Gamma_{i+1}\) 是循环群（这时就说群 \(\Gamma\) 是可解的）。反之，若情况如此，拉格朗日预解式的使用表明 \(K_{i+1}\) 是由向 \(K_i\) 添加某个二项方程的所有根而得到的，从而方程 \(F(x) = 0\) 可用根式求解。\(^{28}\) 于是，“一般”方程——其次数 \(n > 4\)——根式解之不可能，便是如下事实的推论：这个方程的群 \(\Gamma\) 同构于对称群 \(S_n\)，而后者不是可解的。

从 19 世纪中叶起，如我们已经指出的那样（参见第~53~页），代数学家们大大扩展了他们的研究范围，而在此之前，这一范围几乎完全局限于方程的研究。在伽罗瓦发现的启示下，人们看出“根式解”问题只不过是无理数分类这一一般问题的一个相当人为的特殊情形。在整个 19 世纪后半叶，人们将从多个方面进攻的正是后一个问题，许多零散的结果将逐渐积累起来，为施泰尼茨的综合准备了道路。

首先就有代数无理数而言，分类的一条基本原理由伽罗瓦理论提供，它把代数方程的研究化归为其群的研究。事实上，在纯粹代数中，此时研究的主要对象是置换群的理论，对此我们在这里无须谈论（参见第~53~页）。代数数域理论的其他进展来自同一时期数论和代数几何的发展。这些进展主要涉及阐述理论的方式，而且其中大部分应归功于戴德金 {[}79{]}，他引入了域和环的概念，\(^{29}\) 并（联系到他对超复系统的研究）系统地发展了扩张理论的线性方面（{[}79{]}，第 III 卷；第~33~页以下）。也是他把伽罗瓦群视为由所论扩张的自同构所组成，而不再仅仅视为方程诸根的一个置换群；他（就数域）证明了关于自同构线性无关的基本定理（{[}79{]}，第 III 卷，第~29~页）以及伽罗瓦扩张的正规基的存在性（{[}79{]}，第 II 卷，第~433~页）。最后，他着手处理无限次代数扩张的问题，并注意到伽罗瓦理论不能原封不动地应用于它（伽罗瓦群的任意子群并不总是与一个扩张关于某个子扩张的群相同）；而且凭着一种大胆的直觉，他已经开始把伽罗瓦群设想为一个拓扑群 \(^{30}\)——这一思想直到 1928 年才由克鲁尔发展的无限次伽罗瓦扩张理论开花结果 {[}187 d{]}。

与这一演变并行，域上的超越元素的概念变得精确了。超越数的存在性第一次由刘维尔于 1844 年借助一个基于丢番图逼近理论的显式构造得到证明（{[}204 c{]}）；康托尔于 1874 年用关于集合的势的简单考虑给出了另一个“非构造性”证明（{[}47{]}）；最后，埃尔米特于 1873 年证明了 e 的超越性，林德曼于 1882 年用与埃尔米特类似的方法证明了 \(\pi\) 的超越性，从而给化圆为方这一古老问题盖上了最终的封印。\(^{31}\)

至于超越数在代数计算中所起的作用，克罗内克于 1882 年注意到，若 x 在域 K 上是超越的，则域 \(K(x)\) 同构于有理分式域 \(K(X)\)（{[}186 a{]}，第 II 卷，第~253~页）。他还把未定元向域的添加作为其代数数理论阐述的基石（{[}186 a{]}，第 II 卷，第~245-387~页）。另一方面，戴德金和韦伯于同年 {[}80{]} 表明，算术的方法如何能帮助建立代数曲线的理论。由此可以看到算术与代数几何之间的类似在几个方向上显现出来，而这些类似对双方都将显示出极大的丰饶。

在所有这些研究中，出现的域都是由经典数学意义下的“具体”元素——（复）数或复变量的函数——组成的。\(^{32}\) 但早在 1882 年，克罗内克就充分认识到（高斯和伽罗瓦都曾隐约感觉到）这样一个事实：“未定元”在他的理论中所起的作用只是某个代数的基元素的作用，而不是分析意义下的变量的作用（{[}186 a{]}，第 II 卷，第~339-340~页）；1887 年，他发展了这一思想，并将其同一个庞大的纲领联系在一起，该纲领的目标不外乎重新铸造全部数学，摒弃一切不能化归为整数上的代数运算的东西。正是在这个场合，克罗内克接过柯西（{[}56 a{]}，（1），第 X 卷，第~312~页和第~351~页）的一个想法——柯西曾把复数域 C 定义为剩余域 \(\mathbf{R}[X]/(X^{2}+1)\)——并表明代数数理论完全独立于“代数基本定理”，甚至独立于实数理论：所有（有限次的）代数数域都同构于一个剩余域 \(\mathbf{Q}[X]/(f)\)（f 是 \(\mathbf{Q}\) 上的不可约多项式）（{[}186 a{]}，第 III\(_{1}\) 卷，第~211-240~页）。几年之后，H. 韦伯 {[}327 a{]} 在发展域的公理化理论的一份最初草稿时指出，克罗内克的这一方法实际上适用于一切基域 K。韦伯特别指出，可以取 \(\mathbf{Z}/(p)\)（p 为素数）作为 K，从而使“模 p”同余的演算回归到域论；后者诞生于 18 世纪下半叶，与欧拉、拉格朗日、勒让德和高斯的名字相联，而它所显示出的与代数方程理论的类似并未逃过观察者的眼睛；发展这一类似，伽罗瓦（心中惦记着他关于群论的研究）毫不犹豫地引入了模 p 不可约同余的“理想根”，\(^{33}\) 并指出了其主要性质（{[}123{]}，第~113-127~页）。\(^{34}\) 当把克罗内克的方法应用于 \(\mathbf{Z}/(p)\) 时，重新得到的（除术语外）乃是塞尔和戴德金（{[}79{]}，第 I 卷，第~40~页）已经给出的关于这些“伽罗瓦虚数”的理论的表述。

在世纪之交，所有这些“抽象域”的例子之外又添上了若干类型截然不同的新域，即韦罗内塞 {[}318{]} 引入的形式幂级数域，尤其是亨泽尔 {[}157 f{]} 的 p 进域。正是后者的发现促使施泰尼茨（正如他自己明确说明的那样）提炼出所有这些理论共有的抽象观念；他在一部奠基性的著作 {[}294 a{]} 中完成了这项工作，这部著作可以说孕育了代数的现今概念。他系统地推演交换域公理的各种推论，从而引入了素域、可分（代数）元、完备域等概念，定义了扩张的超越次数，并最终证明了任意域的代数闭扩张的存在。

近来，施泰尼茨的理论在若干重要方面得到了补充。一方面，阿廷的工作使伽罗瓦理论的线性特征凸显出来 {[}7 a{]}。另一方面，导子这一普遍概念（仿照经典微分学的形式性质而得）——戴德金（{[}79{]}，第 II 卷，第~412~页）已隐约感觉到它，施泰尼茨在有理分式域的特殊情形下引入了它（{[}294 a{]}，第~209-212~页）——在超越扩张的研究（对现代代数几何至关重要）中得到了成功的运用，尤其是运用于把可分性概念推广到这些扩张 {[}330 d{]}。

\chapter{整除；有序域。}

初等算术运算，尤其是分数的演算，必然导致了对整数整除性的大量经验性验证。但无论是巴比伦人（尽管他们如此精于代数），还是埃及人（尽管他们对分数的计算如此灵巧），似乎都不曾知道支配这些性质的一般法则，这里的开创之功要归于希腊人。他们的算术工作——欧几里得 {[}107{]} 的第 VII 卷和第 IX 卷中有对其精湛的阐述——比起他们在其他数学分支中的那些最美的发现也毫不逊色。两个整数的最大公约数的存在性，在第 VII 卷一开头就已通过以“欧几里得算法”之名著称的程序得到证明；¹ 它是此后一切进展（素数的性质、最小公倍数的存在性与求法等等）的基础；而这一结构的顶冠则是两个出色的定理：其一是证明素数有无穷多个（第 IX 卷，命题 20），其二是给出一个从某些素数出发构造偶完全数的程序（正如欧拉后来所证明的，这个程序实际上给出了所有的偶完全数）。唯有素因子分解的存在性和唯一性没有以一般形式得到证明；不过欧几里得仍明确证明了每个整数都被某个素数整除（第 VII 卷，命题 31），以及下面两个命题（第 IX 卷，命题 13 和 14）：

“如果从一个单位开始，取任意多个数排成常数比的数列{[}即等比数列{]}，且紧接单位的那个数是素数，则最大的那个数除数列中出现的那些数外不能被任何其他数整除”（换句话说，素数的幂 \(p^n\) 只能被 \(p\) 的幂中指数 \(\leq n\) 者整除）。

“如果某数是能被{[}给定的{]}诸素数整除的数中最小的一个，那么除最初{[}给定{]}整除它的那些素数外，它不能被任何其他素数整除”（换句话说，互异素数之积 \(p_1 \ldots p_k\) 除 \(p_1, \ldots, p_k\) 外没有别的素因子）。

因此看来，欧几里得之所以没有陈述一般定理，只是由于他的术语和记法仅足以处理整数的任意幂。\(^{2}\)

尽管细致的研究使人认为，欧几里得的文本中存在着若干先后相继的层次，各自对应于算术发展的一个阶段，\(^{3}\) 但这一演化似乎在公元前 5 世纪初至公元前 4 世纪中叶之间即已完全完成，人们只能赞叹其中展现的敏锐与逻辑的稳健：要再过两千年，我们才能在算术中目睹可比拟的进展。

正是所谓“不定”问题或“丢番图”问题，处于数论此后种种进展的源头。“丢番图方程”一词按今天的用法，从历史上看并非完全有据；一般用它指整系数的代数方程（或方程组），并且只求整数解：当方程是“确定的”、也就是说只有有限多个（实数或复数）解时，这样的问题通常无解；相反，当未知数个数多于方程个数时，它往往有解。不过，尽管丢番图显然是最早考虑“不定”问题的人，他却只是例外地寻求整数解，大多数时候满足于得到一个有理数解 {[}91 a{]}。这类问题他大多能靠代数计算解决，其中未知数的算术性质根本不起作用；\(^{4}\) 因此整除理论在其中只扮演一个很不显眼的角色（“素数”一词只被提及一次（{[}91 a{]}，第 V 卷，问题 9，第 I 卷，第~334-335~页），而互素数的概念也只是在断言两个互素数的商除非其中每一个都是平方数便不可能是平方数的那条定理处才被提到）。\(^{5}\)

不定方程整数解的研究，实际上直到中世纪盛期的中国数学家和印度数学家才真正开始。前者似乎是由于编制历法中的实际问题（其中确定若干天文现象周期的公共周期，恰好构成一个一次的“丢番图”问题）而引到这类思考的；无论如何，联立线性同余方程组的一条解法规则总要归功于他们（无疑在公元 4 世纪至 7 世纪之间）。至于印度人——他们的数学从 5 世纪到 13 世纪处于全盛时期——他们不仅知道如何系统地（通过应用欧几里得算法）处理任意多个未知数的线性丢番图方程组，\(^{6}\) 而且是最早着手并解决了二次问题的人，其中包括“费马方程” \(Nx^{2} + 1 = y^{2}\) 的某些特殊情形（{[}78{]}，第 II 卷，第~87-307~页）。

我们在这里无须追述次数 \textgreater{} 1 的丢番图方程理论的历史；继费马、欧拉、拉格朗日和高斯的工作之后，这一理论在 19 世纪以代数整数理论告终（参见第~93~页以下）。正如我们已经指出的（参见第~58~页），线性方程组的研究在这一时期似乎不再提出值得注意的问题，因而多少受到忽视：特别是，人们既不去寻求任意方程组的解的一般必要条件，也不去描述解的集合。然而，大约在 19 世纪中叶，埃尔米特由于自己在数论方面的研究而用到关于线性丢番图方程的若干引理，尤其是整系数线性代换的一种“简化形”（{[}159{]}，第 I 卷，第~164~页和第~265~页）；最后，在黑格于 1858 年给出了秩等于方程个数的方程组的必要条件之后，H. J. 史密斯于 1861 年以一种“标准形”定义了整元素矩阵的不变因子（{[}287{]}，第 I 卷，第~367-409~页）。

但在此期间，阿贝尔群的概念自高斯引入（参见第~60~页）之后逐渐明确起来，这一概念在数论此后的发展中所占的分量也随之增大。在《算术研究》中对给定判别式的二次型类所成的有限阿贝尔群所作的特别深入的研究里，高斯很快察觉到这些群中有一些并不是循环群：“在这种情形”，他说，“一个基{[}也就是说一个生成元{]}是不够的，必须取两个或更多个，它们通过乘法和合成 \(^{7}\) 将能产生所有的类”（{[}124 a{]}，第 I 卷，第~374-375~页）。高斯说这些话时是否想要描述群分解为循环群的直积，这一点并不确定；不过在《算术研究》的同一页上，他证明了群中存在一个元素，其阶是所有元素的阶的最小公倍数——换句话说，他得到了该群的最大不变因子的存在性（{[}124 a{]}，第 I 卷，第~373~页）——；另一方面，直积的概念他是熟悉的，因为在一篇写于 1801 年但生前未曾发表的手稿中，他勾勒了有限阿贝尔群分解为 \(p\)-群之直积 \(^{8}\) 的一个一般证明（{[}124 a{]}，第 II 卷，第~266~页）。无论如何，1868 年，高斯著作的出版者谢林受这些结果（尤其是他刚刚重新发现的这篇手稿）的启发，证明了（仍是对二次型类群而言）一般分解定理（{[}272{]}，第 I 卷，第~135-148~页），其所用的方法两年后由克罗内克以抽象的方式重新提出（{[}186 a{]}，第 I 卷，第~273-282~页），本质上就是今天仍在使用的方法。至于无挠阿贝尔群，我们已经说过（参见第~64~页以下）椭圆函数与阿贝尔积分的理论——它由高斯、阿贝尔和雅可比发展起来——如何一步步导致认清其结构；把一个无限群分解为单演群之直和的第一个也是最著名的例子，是狄利克雷 1846 年在他关于代数数域的单位的论文中给出的（{[}92{]}，第 I 卷，第~619-644~页）。但直到 1879 年，有限型阿贝尔群的理论与史密斯定理之间的联系才由弗罗贝尼乌斯和施蒂克贝格尔明确地认识到并加以使用（{[}120{]}，§10）。

大约在同一时期，矩阵（具实系数或复系数的）相似理论也趋于完成。线性代换的特征值概念明确出现于常系数线性微分方程组的理论中，拉格朗日把它应用于小振动理论（{[}191{]}，第 I 卷，第~520~页），拉格朗日（{[}191{]}，第 VI 卷，第~655-666~页）和拉普拉斯（{[}193{]}，第 VIII 卷，第~325-366~页）把它应用于行星的“长期”不等式。它也隐含在 18 世纪中叶前后处理的许多其他问题中，比如圆锥曲线或二次曲面的轴的确定（最先由欧拉完成（{[}108 a{]}，（1），第 IX 卷，第~384~页）），又如固体的惯量主轴的研究（也由欧拉加以发展（{[}108 a{]}，（2），第 III 卷，第~200-201~页）；惯量主轴是德·塞格纳于 1755 年发现的）；今天我们知道，正是它（以隐蔽得多的形式）也参与了偏微分方程理论的发端，尤其是弦振动方程。但（且不谈后一情形）这些问题之间的共同渊源直到柯西（{[}56 a{]}，（2），第 V 卷，第~248~页和第 IX 卷，第~174~页）才勉强被认识到。此外，由于这些问题大多涉及对称矩阵，起初主要研究的是这些矩阵的特征值；这里要注意，早在 1826 年，柯西就证明了这些矩阵的特征值在相似下不变，并证明了 3 阶对称矩阵的特征值都是实数（{[}56 a{]}，（2），第 V 卷，第~248~页），三年后他把这一结果推广到任意的实对称矩阵（{[}56 a{]}，（2），第 IX 卷，第~174~页）。⁹ 默比乌斯于 1827 年引入的射影对应的一般概念（{[}223{]}，第 I 卷，第~217~页）很快引出这些变换的分类问题（首先是 2 维和 3 维情形），它不是别的，正是相应矩阵的相似问题；但在很长时间里，这个问题只用 19 世纪中叶风行的“综合”方法来处理，其进展（无论如何相当缓慢）似乎没有对特征值理论产生任何影响。几何中的另一个问题即圆锥曲线束或二次曲面束的分类则不然，从现代观点看，它归结为矩阵 \(U + \lambda V\) 的初等因子的研究，其中 \(U\) 和 \(V\) 是两个对称矩阵；西尔维斯特 1851 年正是沿着这条路线处理这个问题，他仔细考察（目的是为所考虑的束找出“标准形”）矩阵 \(U + \lambda V\) 的各子式在给 \(\lambda\) 代入一个使其行列式为零的值之后变成什么（{[}304{]}，第 I 卷，第~219-240~页）。特征值理论的纯代数方面也在同步进展；于是若干作者（其中有西尔维斯特本人）大约在 1850 年证明了 \(U^n\) 的特征值是 \(U\) 的特征值的 n 次幂；而 1858 年，凯莱在他建立矩阵演算的那篇论文中（{[}58{]}，第 II 卷，第~475-496~页），对任意阶的方阵陈述了“凯莱—哈密顿定理”，\(^{10}\) 他只满足于对 2 阶和 3 阶的矩阵通过直接计算给出证明。最后，1868 年，魏尔斯特拉斯重新采用西尔维斯特的方法，对一个“束” \(U + \lambda V\) 得到了“标准形”，这一次 \(U\) 和 \(V\) 是不必对称的方阵，唯一的限制条件是 det( \(U + \gamma V\) )

不恒为零；他由此给出任意（复元素的）方阵的初等因子的定义，并证明这些初等因子在相似的意义下刻画该方阵（{[}329 a{]}，第 II 卷，第~19-44~页）；这些结果两年后又由若尔当部分地（而且显然是独立地）重新得到 \(^{11}\)（{[}174 a{]}，第~114-125~页）。这里同样是弗罗贝尼乌斯，他于 1879 年表明，魏尔斯特拉斯定理可以容易地从推广到多项式的史密斯理论中导出（{[}119{]}，第 I 卷，第~482-544~页，§13）。

我们刚才提到了单变量多项式的整除理论；多项式除法的问题从代数的一开始就自然会提出，它是乘法运算的逆运算（后者至少对低次多项式而言已为丢番图所知）；但可以想见，在变量的各种幂的一套连贯记法确立之前，很难一般地处理这个问题。事实上，如我们所知的那种多项式“欧几里得”除法的例子，在 17 世纪中叶之前很难找到；\(^{12}\) 而 S. 斯泰芬（他本质上使用了指数记法）似乎是最先想到由此推出“欧几里得算法”的推广、用以求两个多项式的最大公因式的人（{[}295{]}，第 I 卷，第~54-56~页）。有鉴于此，直到 18 世纪中叶，整除性概念一直是有理整数所特有的。是欧拉于 1770 年开启了算术的新的一章，他不无胆量地把整除性概念推广到二次扩张中的整数：为求形如 \(x^{2}+cy^{2}\) 的数（ \(x,y,c\) 为有理整数）的因数，他写下 \(x+y\sqrt{-c}=(p+q\sqrt{-c})(r+s\sqrt{-c})\) （ \(p,q,r,s\) 为有理整数），并对两个成员取范数，从而毫不犹豫地断言，他以这种方式得到了 \(x^{2}+cy^{2}\) 的所有形如 \(p^{2}+cq^{2}\) 的因数（{[}108 a{]}，（1），第 I 卷，第~422~页）。换句话说，欧拉的论证就好像环 \(\mathbb{Z}[\sqrt{-c}]\) 是主理想环一样；稍后，他用一个类似的论证把“无穷递降”方法应用于方程 \(x^{3}+y^{3}=z^{3}\) （这归结为写下 \(p^{2}+3q^{2}\) 是一个立方数，而他通过陈述 \(p+q\sqrt{-3}=(r+s\sqrt{-3})^{3}\) 做到这一点）。但从 1773 年起，拉格朗日证明了（{[}191{]}，第 III 卷，第~695-795~页）形如 \(x^{2}+cy^{2}\) 的数的因数并非总具这种形式，这是那个基本困难的第一个例子；这一困难后来在高斯及其后继者关于单位根域中整除性的研究中，将以清晰得多的面目出现；\(^{13}\) 一般说来，不可能把有理整数的那两条本质的整除性质——最大公因数的存在和素因子分解的唯一性——直接推广到这些域。这里不是详细描述下述情形的地方：库默尔就单位根域 {[}188 b{]}、\(^{14}\) 继而戴德金和克罗内克就任意的代数数域，如何通过创立理想理论而成功地克服这一可怕的障碍——理想理论是现代代数最具决定性的进展之一（参见第~93~页以下）。但戴德金对各种数学理论的基础总是怀有好奇，并不满足于这一成功；他分析了整除关系的机制，在一篇论文中奠定了有序群的现代理论的基础，这篇论文（未曾在其同时代人中激起波澜，并被遗忘了 30 年）按年代算无疑属于关于公理化代数的最早一批工作之列（{[}79, 第 II 卷，第~103-147~页).

从 18 世纪中叶起，寻求“代数基本定理”的一个证明成了当务之急（参见第~74~页）。我们在这里无须回顾达朗贝尔的尝试，他开创了使用无穷小演算的一连串证明（参见第~161~页）。但 1749 年，欧拉以另一种方式处理这个问题（{[}108 a{]}，（1），第 VI 卷，第~78-147~页）：对每个实系数多项式 f，他试图证明分解 \(f = f_{1}f_{2}\) 的存在性，即分解为两个（非常数的）实系数多项式，这将通过对 f 的次数作归纳而给出“基本定理”的一个证明。他甚至指出，如他所言，只需在第一个奇次数的因子处停下即可，因此一切困难都归结为考虑 f 的次数 n 为偶数的情形。于是欧拉只限于研究所求的两个因子都是 n/2 次的情形，他指出，通过一个恰当进行的消元计算，\(f_{1}\) 和 \(f_{2}\) 的未知系数可以用一个实系数方程的根有理地表示出来，该方程的两端项符号相反，因而至少有一个实根。但欧拉的证明只是一个概要，其中若干本质之点被略过未论；直到 1772 年，拉格朗日才成功地解决了这个证明所提出的困难（{[}191{]}，第 III 卷，第~479-516~页），他借助一项极其冗长而极其细致的分析，在使用他（可以说是）新创造的“伽罗瓦”方法上展现出惊人的娴熟（参见第~75~页以下）。

尽管如此，拉格朗日与欧拉及其所有同时代人一样，毫不犹豫地在多项式的一个“根域”中作形式推理（也就是说，用他的语言，考虑该多项式的“虚根”）；他那个时代的数学没有为这类论证提供任何正当性。于是，对 18 世纪那种狂热的形式主义自始便自觉抱有敌意的高斯，在其博士论文中猛烈地奋起反对这种滥用（{[}124 a{]}，第 III 卷，第~3~页）。但如果他没有感觉到这里涉及的是一种表述有缺陷而其本身正确的论证，那他也就不是他自己了。于是我们看到，几年之后（{[}124 a{]}，第 III 卷，第~33~页；另见 {[}124 b{]}），他重新采用欧拉论证的一个较简单的变种——它早在 1759 年就已被德·丰瑟内提出（但后者未能把它圆满完成）——并由此推出“基本定理”的一个新证明，其中他小心避免对“虚”根的任何使用：这些根被未定元的巧妙添加和特殊化所取代。

于是，拓扑学在“基本定理”中所起的作用被归结为单独一条定理：实系数多项式不可能在一个区间内变号而不变为零（多项式情形的波尔查诺定理）。这条定理也是多项式（实系数的）实根的一切分离准则的根基，而这些准则在整个 19 世纪都属于代数最受喜爱的课题之列。\(^{15}\) 在这些研究中，人们不可能不注意到，在那里起本质作用的与其说是 R 的拓扑，不如说是 R 的序结构；\(^{16}\) 例如，多项式情形的波尔查诺定理对全部实代数数所成的域仍然成立。这一思想运动在有序域的抽象理论中达到了它的归宿，该理论由 E. 阿廷和 O. 施赖尔创立（{[}7 b{]} 与 {[}8 a and b{]}）；其更为显著的成果之一，无疑是发现了域上序关系的存在性与该域的纯粹代数性质相联系。
